\properlane{charity-within}{A lamp sustained by charity within}
\subsection{Good deeds whose purpose can fail}
Augustine gives the lamps a wider meaning than Chrysostom does. They signify good works. His examples include feeding the hungry, receiving a stranger, clothing the naked, reconciling opponents, visiting the sick and burying the dead. A person can perform these acts and still belong among the foolish virgins. This is a more searching problem than the failure to perform them: the visible benefit to a neighbor does not by itself disclose why the deed was done.\footnote{Augustine, Sermon 93, §§2--5; NPNF Sermon 43.}

His interpretation begins with the common possessions of the ten. They are all called virgins and all have lamps. Augustine extends virginity to the restraint of the senses from unlawful indulgence and understands the ten as concerning the Church as a whole. Married people are therefore included in his exhortation through fidelity and chastity appropriate to their state. The division within the group cannot be explained by contrasting Christians who have every visible sign of faithfulness with people who have none.

To explain the oil, Augustine turns to Paul's account of charity. Gifts and impressive acts can lack the love that makes them fruitful before God. Augustine's image of oil rising above other liquids illustrates the preeminence of charity. The point is not that works cease to matter; the lamp still has to burn. The question is what sustains it. His examples require the reader to distinguish the deed's visible shape from the inward purpose that accompanies it.

This means that the second interpretation asks something different from the first. Chrysostom addresses the neglected neighbor. Augustine can address the person who already gives generously and is known for it. The same act may be performed out of love of God and neighbor or as a means of securing praise. Observers may correctly report that someone has been fed in either case. They cannot supply the charity missing from the giver's purpose by admiring the result.

The distinction should not make a person abandon useful service until every mixed motive has been eliminated. Augustine tells his hearers to carry oil within while continuing the works represented by the lamps. The practical response is to seek purification of the purpose in doing good. The poor person does not become less entitled to help because the giver discovers vanity. That discovery exposes a further need for conversion within the giver.

\subsection{The vessel and the testimony of conscience}
The vessel is essential to Augustine's explanation. The wise bring something with them, within their own hearts. He calls this the testimony of conscience and inward charity. The foolish rely on the favorable testimony of others. Their works continue to shine while praise is available, but at the judgment each person must answer for his or her own life. The audience that formerly sustained a reputation cannot sustain it there.\footnote{Augustine, Sermon 93, §§9--10.}

Gregory develops the problem by describing practices that can have an admirable exterior: bodily discipline, teaching, generosity to the needy. Any of them can be directed toward human favor. He reads the vessels as hearts in which thoughts are held and contrasts glory kept in conscience with glory sought from other people's mouths. At the arrival, the foolish seek external testimony because they find themselves empty within.\footnote{Gregory the Great, \textit{Homiliae in Evangelia}, 12.1--3.}

Gregory's account gives a particular force to the request for oil. The foolish seek someone to testify to the works that had been seen. Their appeal resembles a demand that companions establish their standing. But the judge knows both deeds and hearts. An accurate list of public actions cannot conceal a purpose known to God. The wise cannot furnish another person's inward testimony by lending their own.

The sellers are therefore flatterers in Augustine and Gregory. This differs materially from Chrysostom's identification of the sellers with the poor. Augustine's sellers provide the praise for which the foolish have lived; their trade has sustained a dependence on approval. The direction to go to them exposes the futility of that dependence at the final encounter. Augustine treats it as a rebuke or mockery rather than a practical instruction that could still repair the failure.\footnote{Augustine, Sermon 93, §§11--14; Gregory, \textit{Homiliae in Evangelia}, 12.3.}

Conscience does not become an independent court whose favorable verdict obliges God. Augustine explicitly interprets the wise virgins' fear of insufficient oil as sober humility. Even a good conscience knows the need to ask forgiveness and acknowledges that God's judgment sees more deeply. Gregory likewise says that a person's own testimony scarcely suffices for that person, much less for another. The inwardness of the preparation therefore excludes both dependence on applause and self-acquittal.

The entrance prayer fits this humility through its request to be heard. The person approaching God seeks acceptance rather than presenting an emotional performance. Augustine's interpretation of Christ's sorrow in Psalm 87 also excludes a mistaken test of inward integrity: distress does not show that the heart lacks grace. A grieving or anxious person may pray truthfully. The relevant examination concerns allegiance and charity, not whether an observer finds the person's manner cheerful.\footnote{Augustine, \textit{Expositions on the Psalms}, Psalm 87.2--4.}

\subsection{The inward gift and the whole person}
Augustine insists that the oil is God's gift. A person can put oil into a vessel but does not create the olive. This image corrects a possible misunderstanding of personal responsibility. Something must truly be present in the person, yet its source is not that person's unaided power. The inward gift can be received, cherished and used without being claimed as a private achievement.\footnote{Augustine, Sermon 93, §§9, 14.}

Bonaventure's Wisdom commentary explains a compatible priority. Wisdom's own light makes her visible, and her approach precedes the seeker's possession of her. The opening of intellect and will requires a response to that approach. Thus the interior life is an attentive reception of what comes from God. It is not sustained by repeatedly inspecting one's spiritual attractiveness or by imagining how one appears to an audience.\footnote{Bonaventure, \textit{Commentarius in librum Sapientiae}, 6:13--17 Vulgate, Quaracchi VI, pp. 147--149.}

The Collect concerns freedom for the service of God. Read alongside this interpretation, its reference to the whole person's service reaches the desire that directs outward action. Freedom from an external obstacle does not by itself free someone from the need to be praised. A person may enjoy ample opportunity and still choose service only when noticed. The prayer can be applied to that inward impediment while retaining its concern for concrete bodily and mental freedom.

The Responsorial Psalm describes desire with soul and flesh, praise with lips and hands, and remembrance on the bed. These details prevent an inward interpretation from reducing faithfulness to an inaccessible private feeling. The body participates in desire and action. Augustine explains that it also participates in the hoped-for resurrection. Remembering God when public attention has ceased is one expression of this enduring desire; the body can rest while the soul remains directed toward God.\footnote{Augustine, \textit{Expositions on the Psalms}, Psalm 62.1--12.}

His conjunction of raised hands and mercy supplies a test of the inward claim. Charity that refuses to act toward a neighbor cannot be established merely by saying that it is invisible. Inward purpose and outward work belong together in his account, though the work alone does not prove the purpose. Aquinas makes the distinction explicit in his exposition of Matthew 25: the virgins expose an inward defect, while the following talents expose negligence of action. He records alternative explanations of the oil and sellers rather than flattening them into one.\footnote{Aquinas, \textit{Super Matthaeum}, chapter 25, lectio 1; the opening of lectio 2 supplies the comparison with the talents, an unappointed contextual passage.}

\subsection{Sleeping without losing charity}
Augustine's explanation of the common sleep becomes especially important here. Sleep cannot mean the cooling of charity if the wise persevere to enter the wedding. Death is shared by the wise and foolish, while the disposition disclosed in the encounter differs. The body ceases its present activity; the value of an act done in charity does not depend on its performer remaining available to receive public praise.\footnote{Augustine, Sermon 93, §§6--9, 17.}

The long second reading promises a bodily gathering of the living and dead with Christ. This gives the life within a bodily future. The promised end is not the survival of a reputation among observers. Aquinas's explanation grounds resurrection in Christ's own rising and allows grief over the bodily separation that death causes. A bereaved person's pain does not invalidate trust in that promise.\footnote{Aquinas, \textit{Super I ad Thessalonicenses}, 4, lectio 2.}

The short form grounds hope in Jesus' death and resurrection; the gathering of the living and dead appears in Paul's continuation. Both forms address persons whose confidence needs strengthening. Neither asks them to establish their worth by displaying composure before others.

In Augustine's exposition, midnight means the unexpected arrival. He criticizes calculations that successively predict an end and are disappointed when time continues. The acclamation likewise gives ignorance of the hour as the reason for watchfulness. The consequence for the inward life is persistent fidelity when no anticipated inspection can be staged. An act whose only reason is the prospect of being seen has exactly the dependence that the vessels expose.\footnote{Augustine, Sermon 93, §8; Matthew 24:42a, 44.}

\subsection{Recognition, receiving and being received}
The Prayer over the Offerings relates the sacramental celebration of the Passion to inward reception. In this interpretation, participation can be examined through the purpose with which one approaches the mystery. An outwardly correct action belongs within a life being converted toward God. Augustine's distinction between deed and charity makes that examination intelligible without implying that the sacrament is merely a display of the participant's intentions.

The shepherd Communion option begins with the Lord who nourishes and leads. Augustine's sacramental explanation concerns Christ's care for the Church and the refreshment of baptism. The gift that sustains the interior life is received within this relation to Christ. Bellarmine's recognition that hardship can persist prevents the shepherd's care from being judged by visible prosperity.\footnote{Augustine, \textit{Expositions on the Psalms}, Psalm 22.1--6; Bellarmine, \textit{Commentary on the Psalms}, O'Sullivan abridgment, Psalm 22.}

The Emmaus option contributes a more active correction. Gregory describes disciples who love and doubt at once, then receive the stranger and come to understand. Their inward condition is tested through the hospitality they actually offer. To withdraw from the neighbor in order to protect an imagined purity of conscience would reverse this sequence. Doing what has been heard opens understanding; the Christ whose words they discuss is received at their table before they recognize him.\footnote{Gregory, \textit{Homiliae in Evangelia}, 23.1--2.}

Gregory then compares this reception with being received by Christ. The connection gives the closed door a personal significance: the final encounter concerns recognition and welcome. Augustine explains the bridegroom's denial of knowledge as refusal or rejection, not an absence of information in God. The request for entry is too late because the time for forming the response has ended. Gregory describes the present as the time in which penitents may still enter through conversion.\footnote{Augustine, Sermon 93, §§15--16; Gregory, \textit{Homiliae in Evangelia}, 12.4--6.}

The Prayer after Communion asks the Spirit to preserve integrity after the gift has been received. That petition is appropriate to a person who discovers how easily a good deed acquires a second purpose. The response is neither despair at mixed motives nor satisfaction with outward correctness. It is renewed dependence on the Spirit, joined to the actual practice of good. The God who sees the heart is also the giver from whom the heart seeks help.

\begin{fourSenses}
\item[Literal.] The ten share visible equipment and sleep; their preparation differs before the arrival makes that difference decisive. Wisdom approaches seekers, the psalm expresses desire in body and soul, and Paul grounds hope in Jesus' death and resurrection. The prayers ask for a truthful and enduring response.
\item[Allegorical.] Christ gathers and judges his Church as Bridegroom. His suffering prayer, shepherding and presence to the disciples disclose the Lord on whom its life depends. The inner gift belongs to his grace, not to a reputation created by other people.
\item[Moral.] Augustine and Gregory direct examination toward the purpose of a good act. The works continue, while charity is sought within them and dependence on praise is corrected. Hospitality and the hands extended in mercy keep this examination answerable to the neighbor.
\item[Anagogical.] Divine recognition discloses the truth of the person beyond public approval. The faithful dead share the promised bodily future with Christ. The final welcome fulfills charity that the Spirit has sustained; a borrowed reputation cannot supply it at the closed door.
\end{fourSenses}
